% !TEX root = ../main.tex
% ---------------------------------------------------------------
% GENERATED FILE - DO NOT EDIT.
% Produced by docs/tools/render_reference.py from the repository source.
% Regenerate with: make docs
% ---------------------------------------------------------------

\apibreak
\section{\texttt{qdrant\_write\_lock}}
\label{func:feedback_intelligence_agent.qdrant_coordination.qdrant_write_lock}\index{Functions!qdrant\_write\_lock}

Lock writers and invalidate cached snapshots before any remote mutation.
\apifield{Usage}
\begin{lstlisting}[style=usage, language=Python]
qdrant_write_lock(lock_path: str | Path) -> Iterator[UUID]
\end{lstlisting}
\apifield{Arguments}
\begin{arglist}
  \item[\texttt{lock\_\allowbreak{}path}] \texttt{str \textbar{} Path}, required.
\end{arglist}
\apifield{Value}\texttt{Iterator[UUID]\allowbreak{}}
\apifield{Details}Use this for every external update to a coordinated collection. The marker is atomically replaced before yielding, so a marker-write failure prevents the remote write and remote failures still invalidate other workers' caches. Hold the context through the complete operation, including cleanup. This is a cooperative local lock, not a distributed lock or Qdrant transaction.
\apifield{Source}\texttt{src/\allowbreak{}feedback\_\allowbreak{}intelligence\_\allowbreak{}agent/\allowbreak{}qdrant\_\allowbreak{}coordination.\allowbreak{}py:41}~(\href{https://github.com/DiogoRibeiro7/feedback-intelligence-agent/blob/46813b067e21ac5d947a7bf539b1e3f81b78149f/src/feedback_intelligence_agent/qdrant_coordination.py#L41}{online}), module \cref{module:feedback_intelligence_agent.qdrant_coordination}
